\section{A positive two-variable representation at all positive orders}
\label{rw:sec:double}

The threshold in Theorem~\ref{subunit:thm:domain} concerns a
\emph{one-variable finite signed measure}. It is not a threshold for
analytic continuation or every positivity property.

\begin{theorem}[Positive double kernel]\label{rw:thm:double}
For every $a,b>0$ and $z\in\realcut$,
\begin{equation}\label{rw:eq:double}
 F_{a,b}(z)=z^2\int_{0<y<x<1}
 \frac{W_{a,b}(x,y)}{(1-xz)(1-yz)}\dd y\dd x,
\end{equation}
where
\begin{equation}\label{rw:eq:W}
 W_{a,b}(x,y)=\frac{(-\log x)^{a-1}(\log(x/y))^{b-1}}
 {\Gamma(a)\Gamma(b)}.
\end{equation}
The weight is strictly positive and has total mass $2^{-a}$.
\end{theorem}
\begin{proof}
For $|z|<1$, expand the two denominators. Substitution $y=xv$
gives, for nonnegative integers $j,k$,
\[
 \int_{0<y<x<1}W_{a,b}(x,y)x^jy^k\dd y\dd x
 =\frac1{(j+k+2)^a(k+1)^b}.
\]
Summing over $j+k=n-2$ yields $c_n$. Absolute convergence justifies
the rearrangement, and taking $j=k=0$ gives the total mass.
On each compact subset of $\realcut$, both denominators are bounded
away from zero uniformly in $x,y\in[0,1]$. The integral is consequently
holomorphic there and supplies the stated continuation.
\end{proof}

\begin{corollary}[Disk nonvanishing below and above the threshold]
\label{rw:cor:disk}
For all $a,b>0$, $F_{a,b}$ has no nonzero zeros in
$\{z:|z|\le1,\ z\ne1\}$. More precisely, when $\Imag z>0$ and
$|z|\le1$,
\begin{equation}\label{rw:eq:diskpick}
 \Imag\frac{F_{a,b}(z)}z>0.
\end{equation}
At any upper-half-disk zero of $\Imag F_{a,b}(z)$, the value
$F_{a,b}(z)$ itself is strictly negative real.
\end{corollary}
\begin{proof}
A direct calculation gives
\[
 \Imag\frac{z}{(1-xz)(1-yz)}
 =\frac{(\Imag z)(1-xy|z|^2)}{|1-xz|^2|1-yz|^2}>0
\]
for $0<y<x<1$, proving \eqref{rw:eq:diskpick}. The lower half follows
by conjugation. On the real interval $z<1$, the integral
$F(z)/z^2$ is strictly positive. Finally, if $F(z)=R\in\R$ with
$\Imag z>0$, then $\Imag(R/z)=-R\Imag z/|z|^2>0$, so $R<0$.
\end{proof}

\begin{corollary}[Gaussian negativity and strict log-convexity]
\label{rw:cor:gaussian}
For every $a,b>0$,
\begin{equation}\label{rw:eq:gaussianpositive}
 -g_{a,b}=\int_{0<y<x<1}W_{a,b}(x,y)
       \frac{x+y}{(1+x^2)(1+y^2)}\dd y\dd x>0,
\end{equation}
and
\begin{equation}\label{rw:eq:gaussianbounds}
 \frac{H_2^{(b)}}{4\,3^a}<-g_{a,b}<\frac{H_2^{(b)}}{3^a}.
\end{equation}
The function $(a,b)\mapsto\Gamma(a)\Gamma(b)(-g_{a,b})$ is strictly
log-convex on the positive quadrant.
\end{corollary}
\begin{proof}
Set $z=\iu$ in \eqref{rw:eq:double} and rationalize the denominator.
Since $1<(1+x^2)(1+y^2)<4$ and
$\int W(x+y)=c_3=H_2^{(b)}/3^a$, the bounds follow.
For log-convexity apply H\"older's inequality to the positive integral
with the gamma factors removed. Its parameter dependence is the
exponential of
$(a-1)\log(-\log x)+(b-1)\log\log(x/y)$.
Equality in H\"older would require a nontrivial real linear combination
of these two logarithms to be constant almost everywhere on the open
triangle, which is impossible. This proves strictness for every
nontrivial line segment of parameter pairs.
\end{proof}

\subsection{All-order integral identities}
\begin{theorem}[Hurwitz--polylogarithm transport]\label{rw:thm:transport}
For $a,b>0$, $b\ne1$, and $z\in\realcut$,
\begin{align}\label{rw:eq:transport}
 F_{a,b}(z)={}&\zeta(b)\Li_a(z)-\frac{\Li_{a+b-1}(z)}{b-1}\notag\\
 &-\frac1{\Gamma(b)}\int_0^\infty t^{b-1}q(t)\Li_a(ze^{-t})\dd t.
\end{align}
At $b=1$ the pole-free identity is
\begin{equation}\label{rw:eq:transport1}
 F_{a,1}(z)=\gamma\Li_a(z)-\partial_a\Li_a(z)
       -\int_0^\infty q(t)\Li_a(ze^{-t})\dd t.
\end{equation}
In particular, the index $a+b-1$ in \eqref{rw:eq:transport} is allowed
to be zero or negative.
\end{theorem}
\begin{proof}
Inside the disk, multiply \eqref{subunit:eq:hurwitz}, with $x=n$, by
$n^{-a}z^n$ and sum. Exponential decay at infinity, boundedness of
$q$, and $b>0$ justify termwise integration. The Hurwitz difference
is $H_{n-1}^{(b)}$, so coefficient comparison gives
\eqref{rw:eq:transport}. At $b=1$, use
$H_{n-1}=\gamma+\log n-\int_0^\infty e^{-nt}q(t)\dd t$;
the sum of $(\log n)n^{-a}z^n$ is $-\partial_a\Li_a(z)$.
The integral expressions on the right are locally holomorphic on the
slit plane. Near $t=0$, $\Li_a(ze^{-t})$ stays bounded on compact
slit sets, and at infinity it is $O(e^{-t})$. Analytic continuation
therefore proves the identities throughout the claimed domain.
Standard conventions for $\Li_s$ are those of DLMF 25.12
\cite{rw:DLMFPolylog}.
\end{proof}

These identities are useful even when no signed Hausdorff measure
exists. They are transformations to a depth-one polylogarithm integral,
not assertions that a depth-two constant has been reduced to a finite
basis of depth-one periods.
